\documentclass[a4paper,11pt]{article}
\usepackage[utf8]{inputenc}
\usepackage[T1]{fontenc}
\usepackage[english]{babel}
\usepackage[left=2.5cm,right=2.5cm,top=2cm,bottom=2cm]{geometry}
\usepackage{hyperref}
\usepackage{xcolor}
\usepackage{microtype}
\usepackage{lmodern}

\definecolor{primary}{RGB}{255, 90, 0}
\hypersetup{colorlinks=true, urlcolor=primary}

\begin{document}
\noindent
\textbf{Lo\"ic Aron MBASSI EWOLO} \\
ENSEA -- Double Dipl\^ome Ing\'enieur \\
\href{mailto:loicaron.mbassiewolo@ensea.fr}{loicaron.mbassiewolo@ensea.fr} \quad (+33) 7 59 52 33 98 \\
\href{https://github.com/Nameless0l}{github.com/Nameless0l} \quad \href{https://mbassiloic.tech}{mbassiloic.tech}

\vspace{1em}
\noindent
Cergy, le \today

\vspace{1em}
\noindent
\textbf{Objet : Candidature -- Ing\'enieur IA, Chef de Projet Technique} \\
\textbf{R\'ef\'erence : Offre d'alternance 12 mois -- SFR}

\vspace{1.5em}
\noindent Madame, Monsieur,

\vspace{0.8em}
\noindent Leader français de la télécommunication en transformation numérique, SFR poursuit des initiatives majeures d'IA pour optimiser réseaux, expérience client et services innovants. Je suis particulièrement intéressé par votre approche de l'IA générative appliquée à la télécom et par les enjeux de direction technique d'équipes cross-fonctionnelles en environnement stratégique. C'est pourquoi je candidats avec enthousiasme pour rejoindre votre équipe en tant qu'Ingénieur IA / Chef de Projet Technique en alternance dès septembre 2026.

\vspace{0.8em}
\noindent J'ai orchestré et diriger un système complet d'agents IA multi-agents sous Google ADK : quatre agents collaboratifs (Planning, Météo, Marché, Recommandations) avec LLM Gemini, résolution automatique de conflits et planning hiérarchique. Cette expérience de direction technique m'a formé à la gestion d'architecture complexe, à la coordination entre composants hétérogènes (LLMs, APIs externes, vector stores), et à la résolution rapide de conflits entre agents autonomes. En parallèle, je coordonne des équipes techniques cross-fonctionnelles (data scientists, backend engineers), gère les roadmaps de features et les priorités budgétaires. Ces responsabilités de leadership technique alimentent directement les compétences de chef de projet que vous recherchez pour vos initiatives IA au sein de SFR.

\vspace{0.8em}
\noindent Depuis avril 2026, je déploie un pipeline NLP complet à l'équipe TRIBE (INRIA Saclay) : fine-tuning BERT en production, tokenization spécialisée, classification multi-labels avec scoring interprétable (0-100). Cette expérience directe de ML/NLP en environnement critique (financement CNRS/ANR, collaboration chercheurs mondiaux) m'a inculqué la rigueur et la rigueur méthodologique indispensables pour diriger des équipes techniques en contexte d'innovation. En outre, ma direction du projet Urban Waste Management (premier prix national CITS contre 75 équipes) m'a formé à la gestion de projets complexes multi-étapes, à l'animation d'équipes hétérogènes, et à la valorisation de résultats techniques auprès d'autres stakeholders. Ces expériences cumulées me positionnent idéalement pour orchestrer les initiatives IA télécommunications chez SFR avec une vision à la fois technique et stratégique.

\vspace{0.8em}
\noindent Je suis disponible dès septembre 2026 pour une alternance 12 mois, ce qui m'permettra de délivrer des initiatives IA stratégiques pour SFR tout en consolidant expertise de leader technique d'équipes. Mes expériences en architecture multi-agents, orchestration complexe, et gestion de projets innovants me permettront d'accélérer vos transformations numériques. Je serais ravi de discuter comment ma vision d'ingénieur IA polyvalent, mes compétences de chef de projet et mon expérience de leadership correspondent aux ambitions de SFR en IA générative télécom. Merci pour cette opportunité exceptionnelle.

\vspace{1.5em}
\noindent Veuillez agr\'eer, Madame, Monsieur, l'expression de mes salutations distingu\'ees.

\vspace{2em}
\noindent \textbf{Lo\"ic Aron MBASSI EWOLO}
\end{document}
